\section{Structure of the three-variable cone}\label{sec:three}

The family in \Cref{family:definition} supplies inequalities missing from the
disjoint-support system. We now identify the part of the nonnegative
quadratic cone in which such inequalities can be needed. The argument also
separates two questions that a hull description must answer: whether the
diagonal caps have been imposed, and whether all inequalities with
nonnegative square coefficients have been found.

\subsection{Closedness and exact extension duality}

The higher-degree cancellations in \eqref{eq:set:disjoint} do not cause a
closure problem. This fact will let us use zero sets to exclude membership
in $\mathcal D_3$ without a separate limiting argument.

\begin{lemma}\label{three:closed}
The cones $\mathcal K_D$ and $\mathcal D_3$ are closed. The quadratic
\[
 q_*(x)=8+4(x_1^2+x_2^2+x_3^2)
\]
belongs to $W\cap\operatorname{int}\mathcal K_D$.
\end{lemma}
\begin{proof}
For a weight $w=w_{A,B}$ and its free coordinates $F$, let
\[
 H_w=\int_{C_3}w(x)v_F(x)v_F(x)^\top\,dx.
\]
This matrix is positive definite: the weight is positive in the interior
of the cube, and a nonzero affine function cannot vanish there. If a
sequence of polynomials in $\mathcal K_D$ converges in coefficients, choose a PSD
Gram matrix $G_{j,w}$ for each weight in each polynomial. Their integrals
are
\[
 \sum_w\tr(H_wG_{j,w}).
\]
These sums are bounded, and every summand is nonnegative. Since
$\tr(H_wG_{j,w})\ge\lambda_{\min}(H_w)\tr(G_{j,w})$, each sequence of Gram
matrices is bounded. A subsequence converges simultaneously in all $27$
blocks to PSD matrices representing the limiting polynomial. Thus $\mathcal K_D$
is closed, and so is its intersection with $W$.

Taking $G_w=I$ in every block gives
\[
 \sum_{A\cap B=\varnothing}w_{A,B}
       \left(1+\sum_{i\in F_{A,B}}x_i^2\right)
       =8+4\sum_i x_i^2.
\]
The linear map from the Gram blocks to $V$ is onto. Indeed, each monomial
in $V$ is a generator: use a constant square for a square-free monomial,
and use $x_i^2$ as the square when its exponent is two. The image of the
interior point consisting of identity Gram matrices is therefore an
interior point of $\mathcal K_D$.
\end{proof}

\begin{proposition}\label{three:extension}
Every linear functional on $W$ that is nonnegative on $\mathcal D_3$ has an
extension to $V$ that is nonnegative on $\mathcal K_D$. The moment relaxation $\mathcal R_D$
is closed, and, for a quadratic $q$,
\begin{equation}\label{eq:three:duality}
 \mathcal L_{m,Y}(q)\ge0\quad\hbox{for every }(m,Y)\in \mathcal R_D
 \quad\Longleftrightarrow\quad q\in \mathcal D_3.
\end{equation}
\end{proposition}
\begin{proof}
Let $T:V^*\to W^*$ be restriction. Because $q_*\in\operatorname{int}\mathcal K_D$,
there is an $\epsilon>0$ such that $q_*+\epsilon u\in \mathcal K_D$ whenever
$\norm{u}\le1$. Hence every $\ell\in \mathcal K_D^*$ satisfies
$\epsilon\norm{\ell}\le\ell(q_*)$. A convergent sequence $T\ell_j$ with
$\ell_j\in \mathcal K_D^*$ has bounded $\ell_j(q_*)$, and therefore a convergent
subsequence of extensions. Consequently $T(\mathcal K_D^*)$ is closed. Its dual is
$\mathcal K_D\cap W=\mathcal D_3$, so the bipolar theorem gives
$T(\mathcal K_D^*)=\mathcal D_3^*$. This proves the extension statement and the closedness of
the mass-one slice $\mathcal R_D$.

To account for normalization in \eqref{eq:three:duality}, note that
$\ell(1)\ge0$ for $\ell\in \mathcal K_D^*$. Positive-mass functionals can be
normalized. If $\ell(1)=0$, add $t$ times evaluation at any cube point,
normalize, and then let $t\downarrow0$. Nonnegativity on the mass-one
slice therefore implies nonnegativity on the entire cone $\mathcal K_D^*$, which
is equivalent to membership in $\mathcal K_D\cap W$.
\end{proof}

The same proof applies to the disjoint-support construction in two
variables. We write $\mathcal D_2$ and $\mathcal R_{D,2}$ for its quadratic cone and moment
relaxation.

\subsection{Diagonal caps and coordinate rounding}

For a coordinate $i$, define the linear map
\begin{equation}\label{eq:three:rounding}
 \rho_i f(x)=(1-x_i)f(x)|_{x_i=0}+x_if(x)|_{x_i=1}.
\end{equation}
For a probability law, composing its moment functional with $\rho_i$
replaces coordinate $i$, conditionally on the original point, by a
Bernoulli variable with the same mean. It keeps first moments and mixed
second moments, and sets the diagonal second moment to the first moment.

\begin{proposition}\label{three:caps}
For every $n$,
\begin{align}
 \mathcal P_n&=\mathcal P_n^+
       +\cone\{x_i(1-x_i):1\le i\le n\},\label{eq:three:cap-cone}\\
 \mathcal Q_n&=\mathcal H_n\cap
       \{(m,Y):Y_{ii}\le m_i\text{ for every }i\}.
       \label{eq:three:cap-hull}
\end{align}
In \eqref{eq:three:cap-cone}, the positive-square-coefficient component can
be chosen to have exactly the same mixed coefficients as the original
quadratic.
\end{proposition}
\begin{proof}
If the coefficient $b_i$ of $x_i^2$ is negative, then
\[
 \rho_iq=q+b_ix_i(1-x_i).
\]
The right side is nonnegative on the cube because it interpolates the
two nonnegative restrictions of $q$. Its $x_i^2$ coefficient is zero,
and no other square or mixed coefficient changes. Applying this operation
to every negative $b_i$ gives
\[
 q=q^++\sum_{b_i<0}(-b_i)x_i(1-x_i),\qquad q^+\in\mathcal P_n^+.
\]
The reverse inclusion is immediate.

For the hull statement, write a point of the right side as a true moment
point plus nonnegative diagonal slack. Take a finite probability law for
the true moment point. Rounding coordinate $i$ changes only its diagonal
moment, from its original value to $m_i$. A mixture of the original and
rounded laws realizes any value in this interval. Thus the assumed cap
allows us to match the prescribed $Y_{ii}$. Treating the coordinates in
turn changes no previously matched moment. This proves membership in
$\mathcal Q_n$. The other inclusion follows from $x_i^2\le x_i$ on the
cube.
\end{proof}

The disjoint-support relaxation has the same rounding property even when
its higher moments do not come from a probability law.

\begin{lemma}\label{three:rounding}
The map $\rho_i$ preserves $V$ and maps $\mathcal K_D$ into $\mathcal K_D$. If
$(m,Y)\in \mathcal R_D$, then both the point obtained by replacing $Y_{ii}$ by
$m_i$ and the point $(m,Y+tE_{ii})$, for $t\ge0$, belong to $\mathcal R_D$.
Consequently,
\begin{equation}\label{eq:three:round-slack}
 \mathcal R_D=\bigl(\mathcal R_D\cap\{Y_{ii}\le m_i\text{ for every }i\}\bigr)
       +\{(0,\diag(s)):s\ge0\}.
\end{equation}
\end{lemma}
\begin{proof}
On monomials in $V$, $\rho_i$ replaces $x_i^2$ by $x_i$ and fixes the
others. For a generator $w_{A,B}L^2$, it fixes the generator if
$i\in A\cup B$. If $i$ is free, it gives the two generators
\[
 w_{A,B\cup\{i\}}\bigl(L|_{x_i=0}\bigr)^2
 +w_{A\cup\{i\},B}\bigl(L|_{x_i=1}\bigr)^2.
\]
Thus $\ell\circ\rho_i$ is feasible whenever $\ell$ is feasible, and its
quadratic moments have the claimed form.

Let $\delta_i$ extract the coefficient of the pure monomial $x_i^2$.
On a generator, this coefficient is zero if $i$ is not free; otherwise it
is $w_{A,B}(0)c_i^2\ge0$, where $c_i$ is the coefficient of $x_i$ in $L$.
Therefore $\ell+t\delta_i$ remains feasible and adds only $t$ to $Y_{ii}$.
Rounding every coordinate for which $Y_{ii}>m_i$, then adding the removed
diagonal slack, proves \eqref{eq:three:round-slack}.
\end{proof}

We use the two-variable exactness theorem of
\citet[Theorem~6]{AnstreicherBurer2010}:
the Shor moment matrix, the box RLT inequalities, and the diagonal caps
describe $\mathcal Q_2$. Here the box RLT inequalities are the moment
forms of products of two bound factors, such as $x_i(1-x_j)\ge0$.

\begin{corollary}\label{three:two-dimensional}
The following statements hold.
\begin{enumerate}[label=(\roman*)]
\item $\mathcal D_2=\mathcal P_2^+$.
\item If $q\in\mathcal P_3^+$ has a zero square coefficient, then
      $q\in \mathcal D_3$.
\item If $q\in\mathcal P_3$ is convex, then $q\in \mathcal D_3$.
\item For every $(m,Y)\in \mathcal R_D$, replacing any one $Y_{ii}$ by $m_i$
      gives a point of $\mathcal H_3$. In particular, a point of $\mathcal R_D$
      outside $\mathcal H_3$ satisfies all three caps strictly.
\end{enumerate}
\end{corollary}
\begin{proof}
Every point of $\mathcal R_{D,2}$ satisfying the caps obeys the Shor and box RLT
conditions, and hence lies in $\mathcal Q_2$. The two-variable version of
\Cref{three:rounding} shows that any other feasible point is such a point
plus nonnegative diagonal slack. Every member of $\mathcal P_2^+$ is
therefore nonnegative on $\mathcal R_{D,2}$; exact extension duality gives (i).

If, for example, $b_3=0$, then
$q=(1-x_3)q|_{x_3=0}+x_3q|_{x_3=1}$. Both restrictions belong to $\mathcal D_2$
by (i). Multiplying their disjoint certificates by $1-x_3$ and $x_3$
produces generators of $\mathcal K_D$, and the resulting sum is the quadratic $q$.
This proves (ii).

For (iii), let $x^*$ minimize $q$ over the cube. Its first-order
optimality conditions express
\[
 \nabla q(x^*)^\top(x-x^*)
  =\sum_{x_i^*=0}\lambda_i x_i+
    \sum_{x_i^*=1}\mu_i(1-x_i),\qquad\lambda_i,\mu_i\ge0.
\]
The Taylor expansion of $q$ is the sum of this expression, the nonnegative
constant $q(x^*)$, and a PSD quadratic form centered at $x^*$. Each term
has a disjoint-support certificate.

For (iv), take any $q\in\mathcal P_3^+$. Its two endpoint restrictions in
coordinate $i$ belong to $\mathcal D_2$, so $\rho_iq\in \mathcal K_D$. A feasible extension
$\ell$ therefore satisfies $(\ell\circ\rho_i)(q)\ge0$. The rounded moments
are nonnegative on all of $\mathcal P_3^+$, and hence belong to
$\mathcal H_3$. If the original point has $Y_{ii}\ge m_i$, it is this
rounded point plus nonnegative diagonal slack, and also belongs to
$\mathcal H_3$.
\end{proof}

\subsection{Exactness for the nonpositive-product sign class}

The mixed coefficients carry a simple invariant under complementation:
the product $c_{12}c_{13}c_{23}$. A nonpositive product means that a cube
symmetry makes the quadratic submodular, that is, makes all its mixed
coefficients nonpositive.

\begin{theorem}\label{three:sign}
If $q\in\mathcal P_3$ and $c_{12}c_{13}c_{23}\le0$, then
\[
 q\in \mathcal D_3+\cone\{x_i(1-x_i):i=1,2,3\}.
\]
If also $b_i\ge0$ for every $i$, then $q\in \mathcal D_3$.
\end{theorem}
\begin{proof}
If one mixed coefficient is zero, independently complement the endpoints
of the other two edges, as necessary, to make those coefficients
nonpositive. If all are nonzero and their product is negative, first make
$c_{12}$ and $c_{13}$ negative by complementing coordinates $2$ and $3$
when needed. The invariant product then makes $c_{23}$ negative as well.
The disjoint cone and the square coefficients are preserved by these
symmetries. We may therefore assume $c_{ij}\le0$ for all $i<j$.

The exactness theorem of \citet[Theorem~1]{BurerNatarajanWillemsen2025}
states that, in at most three variables, a box quadratic with nonpositive
off-diagonal matrix entries has the same minimum as its relaxation
\begin{equation}\label{eq:three:bnw}
 \begin{pmatrix}1&m^\top\\m&Y\end{pmatrix}\succeq0,
 \qquad Y\le m\mathbf1^\top,
\end{equation}
where the latter inequality is entrywise. The symmetric matrix in the
objective has diagonal $b_i$ and off-diagonal entry $c_{ij}/2$.

Suppose first that $b_i\ge0$. Given $(m,Y)\in \mathcal R_D$, use
\eqref{eq:three:round-slack} to write it as $(m,Y')$ satisfying the caps
plus nonnegative diagonal slack. The point $(m,Y')$ is feasible for
\eqref{eq:three:bnw}: its Shor matrix is the unweighted localizing block,
and $Y'_{ij}\le m_i$ for $i\ne j$ follows from the generator
$x_i(1-x_j)$. Exactness gives $\mathcal L_{m,Y'}(q)\ge\min_{C_3}q\ge0$.
Nonnegative square coefficients ensure that adding the diagonal slack
cannot decrease this value. By \Cref{three:extension}, $q\in \mathcal D_3$.

For general square coefficients, use \Cref{three:caps}. Its explicit
decomposition leaves the mixed coefficients unchanged, so the
positive-square-coefficient component belongs to $\mathcal D_3$ by the preceding
argument.
\end{proof}

This theorem concerns the objective of a three-variable box problem. In a
larger constrained model, an inequality on a triple can arise from a dual
combination of the objective and constraints. The mixed coefficients of
the objective alone then do not justify omitting a family block.

\subsection{Extreme rays and the six relevant orientations}

Let $\mathcal F$ be the family in \Cref{family:definition}. For each of its
$24$ distinct cube-symmetry copies, define the closed cone
$\mathcal F_g=\cl\cone\{q\circ g:q\in\mathcal F\}$. Set
\begin{equation}\label{eq:three:combined-cone}
 \mathcal S=\mathcal D_3+\sum_g\mathcal F_g,
 \qquad
 \mathcal R=\mathcal R_D\cap\{(m,Y):\mathcal L_{m,Y}(f)\ge0\text{ for all }f\in \mathcal F_g, g\}.
\end{equation}
By \Cref{family:lmi}, the additional inequalities defining $\mathcal R$ have a
finite SDP representation.

\begin{lemma}\label{three:sum}
The cone $\mathcal S$ is closed, and $\mathcal S$ is precisely the cone of quadratics
nonnegative on $\mathcal R$. An extreme ray of $\mathcal P_3^+$ belongs to $\mathcal S$ if
and only if it belongs to $\mathcal D_3$ or is a limit of normalized members of one
family copy.
\end{lemma}
\begin{proof}
The functional $J(q)=\int_{C_3}q(x)\,dx$ is strictly positive on every
nonzero element of $\mathcal P_3$. The base
$\{q\in\mathcal P_3:J(q)=1\}$ is compact: otherwise a sequence of
unbounded norms, normalized in coefficient norm, would have a nonzero
nonnegative limit with integral zero. Each summand in $\mathcal S$ is a closed
subcone of $\mathcal P_3^+$. For a convergent sequence of sums, its total
integral bounds the integral, and thus the norm, of every summand. Taking
subsequences proves closedness of $\mathcal S$.

In homogeneous coordinates the moment cone defining $\mathcal R$ is the
intersection of $\mathcal D_3^*$ with the $\mathcal F_g^*$. Its dual is the closure of the
sum of their duals, which is $\mathcal S$ by closedness. The normalization argument
in \Cref{three:extension} applies here too: evaluation at a cube point is
feasible for all the cones. This proves the asserted duality.

If an extreme ray of $\mathcal P_3^+$ is a sum of elements from its
subcones, every nonzero summand lies on that ray. The unit-integral base
of $\mathcal F_g$ is the convex hull of the compact set of limits of
unit-integral members of its family copy. An extreme point of the larger
base of $\mathcal P_3^+$ cannot be a nontrivial convex combination of
these limits. This proves the characterization.
\end{proof}

\begin{corollary}\label{three:reduction}
Every $q\in\mathcal P_3^+\setminus \mathcal D_3$ has strictly positive square
coefficients, an indefinite Hessian, and
$c_{12}c_{13}c_{23}>0$. After complementing at most one coordinate, its
three mixed coefficients are strictly positive. Among the $24$ family
copies, only six can approach an extreme ray with this strict sign
pattern: for each choice of the distinguished coordinate $z$, complement
either $z$ or both of the other coordinates.
\end{corollary}
\begin{proof}
The first claims follow from \Cref{three:two-dimensional,three:sign}.
A positive product means that either all coefficients are positive or
exactly two are negative. In the latter case, complement their common
coordinate.

For the original family,
\[
 c_{xz}=-2d_3(d_1+k)\le0,\qquad
 c_{yz}=-2d_3(d_2+k)\le0.
\]
To obtain three strictly positive coefficients after complementation,
the original $c_{xy}$ must be positive. The pattern $(+,-,-)$ becomes
$(+,+,+)$ exactly when the complemented set is $\{z\}$ or $\{x,y\}$.
There are three choices of $z$. A sequence approaching a strictly positive
pattern eventually has that same strict pattern, so no other orientation
can contribute such a limit.
\end{proof}

An extreme ray of $\mathcal P_3^+$ with all square coefficients strictly
positive is also extreme in $\mathcal P_3$. To see this, let
$q=p_1+p_2$ with $p_j\in\mathcal P_3$. For sufficiently small
$\epsilon>0$, both $q+\epsilon(p_1-p_2)$ and
$q-\epsilon(p_1-p_2)$ remain in $\mathcal P_3^+$: they are positive
linear combinations of $p_1,p_2$, and their square coefficients remain
positive. Extremality in the smaller cone forces $p_1-p_2$ to be
proportional to $q$, and hence both summands lie on its ray.

The reduction is a restriction on possible missing inequalities; it does
not prove that the six orientations suffice. The complete statement
still to be decided is
\begin{equation}\label{eq:three:completeness}
 \mathcal P_3^+=\mathcal S.
\end{equation}
It is equivalent to $\mathcal H_3=\mathcal R$, and also to
\begin{equation}\label{eq:three:full-completeness}
 \mathcal Q_3=\mathcal R\cap\{Y_{ii}\le m_i\text{ for every }i\}.
\end{equation}
Indeed, one implication follows from \Cref{three:caps}. Conversely, if
\eqref{eq:three:full-completeness} holds, a point of $\mathcal R$ either has some
$Y_{ii}\ge m_i$, and is in $\mathcal H_3$ by
\Cref{three:two-dimensional}, or satisfies every cap and is in
$\mathcal Q_3$. Thus $\mathcal R=\mathcal H_3$. Coordinate rounding maps $\mathcal R$ into
$\mathcal H_3\subseteq \mathcal R$ as well.

To prove \eqref{eq:three:completeness}, it is enough to treat extreme rays
with all mixed coefficients strictly positive. The compact-base argument
in \Cref{three:sum}, together with the finite-dimensional
Krein--Milman theorem, reduces cone equality to its extreme rays, and
\Cref{three:reduction} removes all other cases. The computational evidence
in \Cref{comp:section} supports this equality but does not establish it.

\subsection{Zero-set obstructions and boundary extreme rays}

Zero sets explain why weighted affine squares can fail. For a vertex $v$
and a set of free coordinates $F$, put
\[
 U_F(v)=\{x\in C_3:x_j\ne1-v_j\text{ for every }j\notin F\},
\]
and let $\pi_F$ be projection onto $F$.

\begin{lemma}\label{three:blocking}
Let $q\in\mathcal P_3$, with zero set $Z$, and suppose $q(v)>0$ at a
vertex $v$. If
\[
 \pi_F(v)\in\operatorname{aff}\bigl(\pi_F(Z\cap U_F(v))\bigr)
 \qquad\text{for every }F\subseteq\{1,2,3\},
\]
then $q\notin \mathcal D_3$. For $F=\varnothing$, the condition means that
$Z\cap U_\varnothing(v)$ is nonempty.
\end{lemma}
\begin{proof}
In a disjoint certificate for $q$, every nonnegative summand vanishes on
$Z$. Since $q(v)>0$, some summand $w_{A,B}L^2$ has positive weight at $v$
and $L(v)\ne0$. The positive locus of this weight is $U_F(v)$ for its
free coordinates $F$. Hence $L$ vanishes on the projected zero set in the
display, and on its affine hull. The hypothesis forces $L(v)=0$, a
contradiction.
\end{proof}

Because $\mathcal D_3$ is closed, this criterion excludes limiting disjoint
certificates as well as exact ones. By contrast, three interior edge zeros
incident to one vertex cannot yield a missing ray.

\begin{lemma}\label{three:three-edges}
Suppose $q\in\mathcal P_3$, $q(0)>0$, and
$q(\alpha_i e_i)=0$ with $0<\alpha_i<1$ for $i=1,2,3$. Then
\[
 q(x)=q(0)\left[
       \left(1-\sum_i\frac{x_i}{\alpha_i}\right)^2
       +2\sum_{i<j}r_{ij}\frac{x_ix_j}{\alpha_i\alpha_j}\right],
 \qquad r_{ij}\ge0.
\]
In particular, $q\in \mathcal D_3$; if it is extreme in $\mathcal P_3$, it is an
affine square.
\end{lemma}
\begin{proof}
The nonnegative restriction to edge $i$ has an interior double zero, and
equals $q(0)(1-x_i/\alpha_i)^2$. This fixes every coefficient except the
three mixed coefficients and gives the stated representation with real
$r_{ij}$. Evaluation at
$(\alpha_i e_i+\alpha_j e_j)/2$ gives $q(0)r_{ij}/2\ge0$. The displayed
terms have disjoint certificates. If some $r_{ij}>0$, the decomposition
contains the nonproportional nonnegative polynomials $x_ix_j$ and the
affine square, and the ray is not extreme.
\end{proof}

The exposed family rays of \Cref{family:exposed} persist at several
parameter boundaries. Some edge zeros reach a vertex; their tangency
conditions then matter as much as their values.

\begin{theorem}\label{three:boundary}
Let $k>0$ and $d_1,d_2\ge0$. A family member
$q_{h,d,k}$ spans an extreme ray of $\mathcal P_3$, and hence of
$\mathcal P_3^+$, in each of the following regimes:
\begin{enumerate}[label=(\roman*)]
\item $0<h<\min(d_1,d_2)$ and $d_3=d_1+d_2-h+k$;
\item $h=0$ and $d_3\ge d_1+d_2+k$;
\item $h=d_1<d_2$ and $d_3\ge d_2+k$, or the same regime with
      $d_1,d_2$ exchanged.
\end{enumerate}
The assertions include zero parameters $d_1$ or $d_2$ whenever their
regime allows them.
\end{theorem}
\begin{proof}
The family is nonnegative by \Cref{family:valid}. Its edge restrictions
and the vertex tangencies described in \Cref{app:three-boundary} give
double zeros, including those at edge endpoints. In any decomposition
into nonnegative quadratics, both summands inherit every zero and every
listed tangent edge derivative. At an interior edge zero this follows
from ordinary first-order optimality. At a vertex, their one-sided
derivatives into the edge are nonnegative and sum to zero, so both vanish.
The coefficient reconstruction in \Cref{app:three-boundary} shows that
these linear contact conditions have a one-dimensional kernel in each
regime. Thus both summands are multiples of $q_{h,d,k}$.
\end{proof}

In regime (ii), a zero parameter produces a whole zero edge, for example
$q(x,0,0)=d_1^2x^2$ when $h=0$. The proof does not require isolated zeros.
Such members have a zero square coefficient and belong to $\mathcal D_3$ by
\Cref{three:two-dimensional}. Extremality therefore does not by itself
imply that a family member strengthens the disjoint-support system.
